\chapter{คู่มือการใช้งานระบบ}

% Text for the first appendix goes here.



% Text for a section in the first appendix goes here.

% test ทดสอบฟอนต์ serif ภาษาไทย

% \textsf{test ทดสอบฟอนต์ sans serif ภาษาไทย}

% \verb+test ทดสอบฟอนต์ teletype ภาษาไทย+

% \texttt{test ทดสอบฟอนต์ teletype ภาษาไทย}

% \textbf{ตัวหนา serif ภาษาไทย \textsf{sans serif ภาษาไทย} \texttt{teletype ภาษาไทย}}

% \textit{ตัวเอียง serif ภาษาไทย \textsf{sans serif ภาษาไทย} \texttt{teletype ภาษาไทย}}

% \textbf{\textit{ตัวหนาเอียง serif ภาษาไทย \textsf{sans serif ภาษาไทย} \texttt{teletype ภาษาไทย}}}

% \section{Appendix section}
ในส่วนนี้จะเป็นคู่มือการใช้งานระบบ Solar Watch ซึ่งเป็นระบบสำหรับการดูข้อมูลสภาพอากาศที่ถูกส่งมากจากสถานีวัดสภาพอากาศขนาดย่อมและการแสดงผลข้อมูลการพยากรณ์ปริมาณการผลิตไฟฟ้าจากโซลาร์เซลล์ล่วงหน้า 24 ชั่วโมง โดยผู้ใช้งานสามารถเข้าถึงระบบได้ผ่านทางเว็บไซต์ที และสามารถดูข้อมูลได้ตามความต้องการของผู้ใช้งาน
% \url{https://solarwatch.cpe.eng.cmu.ac.th/}

\section{วิธีการใช้งานระบบ Solar Watch}

ผู้ใช้เข้าสู่เว็บไซต์ของระบบ Solar Watch ที่ \url{https://solarwatch.cpe.eng.cmu.ac.th/}
% \begin{center}
%     \includegraphics[width=\textwidth]{Home.jpg}
%     \caption{}
%     \label{fig:Home}
% \end{center}

    \begin{figure}[htbp]
        \centering
        \includegraphics[width=\textwidth]{Home.jpg} 
        \caption{หน้าจอหลักเมื่อผู้ใช้เข้าสู่เว็บไซต์ Solar Watch}
        \label{fig:Home}
    \end{figure}
เมื่อเปิดหน้าเเว็บไซต์ขึ้นมา ผู้ใช้จะเห็นหน้าจอหลักของระบบ Solar Watch ซึ่งประกอบไปด้วยเมนูต่าง ๆ ที่ผู้ใช้สามารถเลือกใช้งานได้ตามความต้องการ โดยเมนูหลักจะประกอบไปด้วย
\begin{itemize}
    \item เมนู Station: เป็นเมนูที่จะแสดงข้อมูลสถานีวัดสภาพอากาศขนาดย่อมและข้อมูลการพยากรณ์ปริมาณการผลิตไฟฟ้าจากโซลาร์เซลล์ล่วงหน้า 24 ชั่วโมง โดยผู้ใช้สามารถเลือกดูข้อมูลของแต่ละสถานีได้ตามต้องการ (ในปุจจุบันมีสถานีวัดสภาพอากาศขนาดย่อมทั้งหมด 1 สถานี ได้แก่ 30 years building ซึ่งตั้งอยู่ที่ ตึก 30 ปี คณะวิศวกรรมศาสตร์ มหาวิทยาลัยเชียงใหม่) 
    \item เมนู Setting Station: เป็นเมนูทีสำหรับเจ้าหน้าที่สำหรับการเช็คสถานะการทำงานของสถานีวัดสภาพอากาศขนาดย่อมและสถานะการทำงานของเซนเซอร์ตที่ติดตั้งอยู๋ที่สถานีวัดสภาพอากาศขนาดย่อม โดยเจ้าหน้าที่สามารถเลือกดูข้อมูลของแต่ละสถานีได้ตามต้องการ (ในปุจจุบันมีสถานีวัดสภาพอากาศขนาดย่อมทั้งหมด 1 สถานี ได้แก่ 30 years building ซึ่งตั้งอยู่ที่ ตึก 30 ปี คณะวิศวกรรมศาสตร์ มหาวิทยาลัยเชียงใหม่)
\end{itemize}

\begin{enumerate}
\item สำหรับผู้ใช้ทั่วไปที่ต้องการดูข้อมูลสภาพอากาศและการพยากรณ์ปริมาณการผลิตไฟฟ้าจากโซลาร์เซลล์ล่วงหน้า 24 ชั่วโมง ให้ผู้ใช้คลิกที่เมนู Station เพื่อเข้าสู่หน้าจอแสดงข้อมูลสภาพอากาศและการพยากรณ์ปริมาณการผลิตไฟฟ้าจากโซลาร์เซลล์ล่วงหน้า 24 ชั่วโมง
\begin{itemize}
    \item[]1.1 เมื่อผู้ใช้คลิกปุ่ม Station แล้วจะมีรายการสถานีวัดสภาพอากาศขนาดย่อมที่มีอยู่ในระบบแสดงขึ้นมา โดยผู้ใช้สามารถเลือกสถานีที่ต้องการดูข้อมูลได้ตามต้องการ (ในปุจจุบันมีสถานีวัดสภาพอากาศขนาดย่อมทั้งหมด 1 สถานี ได้แก่ 30 years building ซึ่งตั้งอยู่ที่ ตึก 30 ปี คณะวิศวกรรมศาสตร์ มหาวิทยาลัยเชียงใหม่)
    \begin{figure}[H]
            \centering
            \includegraphics[width=\textwidth]{select_station.jpg} 
            \caption{รายการสถานีวัดสภาพอากาศขนาดย่อมที่มีอยู่ในระบบ}
            \label{fig:select_station}
    \end{figure}

    \item[]1.2 เมื่อผู้ใช้เลือกสถานีที่ต้องการดูข้อมูลแล้ว ผู้ใช้จะเห็นหน้าจอที่แสดงข้อมูลสภาพอากาศที่ถูกส่งมากจากสถานีวัดสภาพอากาศขนาดย่อมและข้อมูลการพยากรณ์ปริมาณการผลิตไฟฟ้าจากโซลาร์เซลล์ล่วงหน้า 24 ชั่วโมง 
    \begin{figure}[H]
            \centering
            \includegraphics[width=\textwidth]{predicgraph.png} 
            \caption{หน้าจอแสดงข้อมูลสภาพอากาศและกราฟการพยากรณ์ปริมาณการผลิตไฟฟ้าจากโซลาร์เซลล์ล่วงหน้า 24 ชั่วโมง}
            \label{fig:predicgraph}
    \end{figure}

    \item[]1.3 เมนูด้านซ้ายมือของหน้าจอแสดงข้อมูลสภาพอากาศและการพยากรณ์ปริมาณการผลิตไฟฟ้าจากโซลาร์เซลล์ล่วงหน้า 24 ชั่วโมง จะมีเมนูสำหรับเลือกดูข้อมูลสภาพอากาศให้ผู้ใช้สามารถเลือกดูได้โดยจะแสดงข้อมูลสภาพอากาศรายชั่วโมงของวันที่ที่ผู้ใช้เลือกโดยเมนูจะมีเมนูแสดงข้อมูล ความเข้มแสงและอุณหภูมิแผงโซลาร์เซลล์ อุณหภูมิและความชื้นสัมพัทธ์ ความกดอากาศ ความเร็วลม และทิศทางลม
    \begin{figure}[H]
            \centering
            \includegraphics[width=\textwidth]{selectmenu.png} 
            \caption{เมนูสำหรับเลือกดูข้อมูลสภาพอากาศรายชั่วโมง}
            \label{fig:selectmenu}
    \end{figure}

    \item[]1.4 ผู้ใช้สามารถเลือกวันที่ที่ต้องการดูข้อมูลได้โดยการคลิกปุ่มวันที่ที่อยู่ด้านบนกราฟทางขวามือของหน้าจอ
    \begin{figure}[H]
            \centering
            \includegraphics[width=\textwidth]{date.png} 
            \caption{ปุ่มสำหรับเลือกวันที่ที่ต้องการดูข้อมูล}
            \label{fig:date}
    \end{figure}

    \item[]1.5 เมื่อผู้ใช้เลือกวันที่ที่ต้องการดูข้อมูลแล้ว ผู้ใช้จะเห็นข้อมูลสภาพอากาศวันที่ที่ผู้ใช้เลือก โดยเมื่อผู้ใช้เลือกวันที่ย้อนหลังเมนู Electric Current กราฟจะแสดงข้อมูลปริมาณการผลิตไฟฟ้าจากแผงโซลาร์เซลล์ที่ผลิตได้จริง
    \begin{figure}[H]
            \centering
            \includegraphics[width=\textwidth]{actualgraph.png} 
            \caption{หน้าจอแสดงข้อมูลสภาพอากาศเฉลี่ยและกราฟแสดงปริมาณการผลิตไฟฟ้าจากแผงโซลาร์เซลล์ที่ผลิตได้จริง}
            \label{fig:actualgraph}
    \end{figure}

    \item[]1.6 ตัวอย่างเมนูสำหรับเลือกดูข้อมูลสภาพอากาศรายชั่วโมงของวันที่ที่ผู้ใช้เลือก
    \begin{figure}[H]
            \centering
            \includegraphics[width=\textwidth]{sunlightandsolar.png} 
            \caption{ข้อมูลสภาพอากาศรายชั่วโมงของวันที่ที่ผู้ใช้เลือก โดยแสดงข้อมูลความเข้มแสงและอุณหภูมิแผงโซลาร์เซลล์}
            \label{fig:sunlightandsolar}
    \end{figure}
    \begin{figure}[H]
            \centering
            \includegraphics[width=\textwidth]{tempandhum.png} 
            \caption{ข้อมูลสภาพอากาศรายชั่วโมงของวันที่ที่ผู้ใช้เลือก โดยแสดงข้อมูลอุณหภูมิและความชื้นสัมพัทธ์}
            \label{fig:tempandhum}
    \end{figure}
    \begin{figure}[H]
            \centering
            \includegraphics[width=\textwidth]{pressure.jpg} 
            \caption{ข้อมูลสภาพอากาศรายชั่วโมงของวันที่ที่ผู้ใช้เลือก โดยแสดงข้อมูลความกดอากาศ}
            \label{fig:pressure}
    \end{figure}
    \begin{figure}[H]
            \centering
            \includegraphics[width=\textwidth]{winddirection.jpg} 
            \caption{ข้อมูลสภาพอากาศรายชั่วโมงของวันที่ที่ผู้ใช้เลือก โดยแสดงข้อมูลทิศทางลม}
            \label{fig:winddirection}
    \end{figure}

\end{itemize}

\item สำหรับเจ้าหน้าที่ที่ต้องการเช็คสถานะการทำงานของสถานีวัดสภาพอากาศขนาดย่อมและสถานะการทำงานของเซนเซอร์ตที่ติดตั้งอยู๋ที่สถานีวัดสภาพอากาศขนาดย่อม ให้เจ้าหน้าที่คลิกที่เมนู Setting Station เพื่อเข้าสู่หน้าจอแสดงข้อมูลสถานะการทำงานของสถานีวัดสภาพอากาศขนาดย่อมและสถานะการทำงานของเซนเซอร์ตที่ติดตั้งอยู๋ที่สถานีวัดสภาพอากาศขนาดย่อม
\begin{itemize}
    \item[]2.1 เมื่อเจ้าหน้าที่คลิกปุ่ม Setting Stationจะต้องทำการกรอกชื่อผู้ใช้งานและรหัสผ่านเพื่อเข้าสู่ระบบ
    \begin{figure}[H]
            \centering
            
\includegraphics[width=\textwidth]{login.jpg} 
            \caption{หน้าจอสำหรับกรอกชื่อผู้ใช้งานและรหัสผ่านเพื่อเข้าสู่ระบบ}
            \label{fig:login}
    \end{figure}

    \item[]2.2 เมื่อเจ้าหน้าที่กรอกชื่อผู้ใช้งานและรหัสผ่านถูกต้องแล้วจะเข้าสู่หน้าจอแสดงข้อมูลสถานะการทำงานสถานะการทำงานของเซนเซอร์ตที่ติดตั้งอยู๋ที่สถานีวัดสภาพอากาศขนาดย่อม โดยเจ้าหน้าที่สามารถเลือกดูข้อมูลของแต่ละสถานีได้ตามต้องการ (ในปุจจุบันมีสถานีวัดสภาพอากาศขนาดย่อมทั้งหมด 1 สถานี ได้แก่ 30 years building ซึ่งตั้งอยู่ที่ ตึก 30 ปี คณะวิศวกรรมศาสตร์ มหาวิทยาลัยเชียงใหม่) โดยการแสดงสถานะการทำงานจะมีดังนี้
    \begin{itemize}
        \item สถานีวัดข้อมูลสภาพอากาศขนาดย่อมทำงานปกติ: สถานะของเซนเซอร์จะแสดงผลเป็น Online ทั้งหมดและไม่มีการแสดงผลหน้าจอสีแดง
        \begin{figure}[H]
            \centering
            \includegraphics[width=\textwidth]{statusstationOK.png} 
            \caption{สถานะการทำงานของสถานีวัดสภาพอากาศขนาดย่อมทำงานปกติ}
            \label{fig:statusstationOK}
        \end{figure}

        \item สถานีวัดข้อมูลสภาพอากาศขนาดย่อมทำงานผิดปกติ: สถานะของเซนเซอร์จะแสดงผลเป็น Offline อย่างน้อย 1 เซนเซอร์หรือหากไม่มีการส่งข้อมูลมาทุกๆ 15 นาทีหน้าจอจะมีการแสดงผลเป็นสีแดงเพื่อแจ้งเตือนว่ามีปัญหาการทำงานของสถานีวัดสภาพอากาศขนาดย่อม
        \begin{figure}[H]
            \centering
            \includegraphics[width=\textwidth]{statusstationNG.png} 
            \caption{สถานะการทำงานของสถานีวัดสภาพอากาศขนาดย่อมทำงานผิดปกติ}
            \label{fig:statusstationNG}
        \end{figure}
    \end{itemize}

\end{itemize}

\end{enumerate}

% \chapter{\ifenglish Manual\else คู่มือการใช้งานระบบ\fi}

% Manual goes here.
